\documentclass{article}
\usepackage{amsmath}
\usepackage{mathtools}
\begin{document}

\DeclarePairedDelimiter\bra{\langle}{\rvert}
\DeclarePairedDelimiter\ket{\lvert}{\rangle}
\DeclarePairedDelimiterX\braket[2]{\langle}{\rangle}{#1\,\delimsize\vert\,\mathopen{}#2}

Dans ce papier, l'attention sera essentiellement tourné vers les matrices et le patrice de rotation dans le cas continu.
Soit \( \ket{\psi} \), le vecteur de Jones représentant la direction du champ électrique.

Soit \( R(\theta) \), la matrice de la rotation du champ électrique par la polarization linéaire.
Cette matrice peut être écrite comme:
$$
R(\theta) = 
\begin{pmatrix}
	\cos(\theta) & \frac{d}{d\theta}\cos(\theta) \\
	\sin(\theta) & \frac{d}{d\theta}\sin(\theta)
\end{pmatrix}
=
\begin{pmatrix}
	\cos(\theta) & -\sin(\theta) \\
	\sin(\theta) & \cos(\theta)
\end{pmatrix}
$$

Nous allons maintenant prendre une valeur \( d\theta \to 0 \).
$$
R(d\theta) = 
\begin{pmatrix}
	\cos(d\theta) & -\sin(d\theta) \\
	\sin(d\theta) & \cos(d\theta)
\end{pmatrix}
=_{d\theta \to 0}
\begin{pmatrix}
	1 & -d\theta \\
	d\theta & 1
\end{pmatrix}
=
\begin{pmatrix}
	0 & -d\theta \\
	d\theta & 0
\end{pmatrix}
+ I
$$
$$
R(d\theta) = d\theta
\begin{pmatrix}
	0 & -1 \\
	1 & 0
\end{pmatrix}
+ I
= d\theta J + I
$$
On remarque que \( J^2 = -I \), ou plus généralement que:
$$
J^n =
\left\{
\begin{array}{ll}
	I & \quad n \equiv 0 \pmod{4} \\
	J & \quad n \equiv 1 \pmod{4} \\
	-I & \quad n \equiv 2 \pmod{4} \\
	-J & \quad n \equiv 3 \pmod{4}
\end{array}
\right.
$$

\( d\theta \) est une partie infinétisémal de \( \theta \), donc:

$$ 
\lim_{N \to \infty} \lim_{d\theta \to 0} R(d\theta)^N = R(\theta)
$$
$$
R(d\theta)^N = (J\theta + I)^N = e^{J\theta}
$$
Donc la matrice de la rotation continue se traduit comme une exponentielle de matrice. Avec les séries de Taylor, cette exponentielle s'ecrit comme étant:
$$
\sum_{n = 0}^{\infty} \frac{J^n \theta^n}{n!} = I + J\theta - I \frac{\theta^2}{2!} - J \frac{\theta^3}{3!} + I \frac{\theta^4}{4!} ...
$$
$$
\sum_{n = 0}^{\infty} \frac{J^n \theta^n}{n!} = I \sum_{n = 0}^{\infty} \frac{(-1)^n \theta^{2n}}{(2n)!} + J \sum_{n = 0}^{\infty} \frac{(-1)^n \theta^{2n+1}}{(2n+1)!}
$$
On peut reconnaître les séries de Taylor pour sin et cos.
Donc cela nous amène à la formule d'Euler pour les matrices:
$$
e^{J\theta} = I \cos(\theta) + J \sin(\theta)
$$
Nous pouvons de cette maniére exprimer la rotation totale que subit la lumière lorsqu'il y a des polariseurs en continu avec des fonctions trigonométriques.

\end{document}
